\section{CMU PhD：工程与科学的边界}

\subsection{从大疆到学术界的转型动机}

2018年，杨硕在大疆飞控系统已经非常完善的背景下，选择前往CMU攻读机器人学博士。他的核心动机是对人形机器人等复杂机器人系统的长期信念。当时人形机器人在科研上还是一个未被很好解决的问题，杨硕想要深入研究它。

\begin{warningbox}{对强化学习的早期误判}
杨硕坦言，2018年入学时他对强化学习持比较悲观的态度，选择了CMU而非Berkeley，专注于模型预测控制（MPC）等传统方法。回看7年后的技术发展，他承认"当时选择肯定是不对的"，如果能重来，可能会去Berkeley直接研究强化学习，这对现在的工作会更有帮助。他在模型训练的很多细节上掌握得"没有那么深入"。不过，他也指出MPC虽然不会是主流，但其中的很多思想会继续存活，被应用在强化学习中。
\end{warningbox}

\subsection{冯·卡门的定义：科学家与工程师}

杨硕引用了钱学森的导师冯·卡门（Theodore von K\'arm\'an）的经典论述来阐明科学与工程的区别：

\begin{quote}
Scientist 是解释已有的东西，Engineer 是创造不存在的东西。
\end{quote}

这一认知是杨硕花了很长时间才理解清楚的。在读PhD之前，他以工程思维认为"四旋翼已经被研究透了，不值得再研究"。但实际上，PhD的工作本质是\textbf{解释科学}——例如为什么方法B比方法A效果好？是网络结构、训练方法还是其他因素？四旋翼作为一个尺寸小、飞行灵活、工程问题已基本解决的平台，恰恰是验证导航和路径规划等算法的理想科研工具。

\begin{knowledgebox}{杨硕关于"研究四旋翼"判断的自我纠正}
杨硕在2018年的知乎文章中建议学生不要再研究四旋翼飞行器，因为"已经被研究透了"。他后来承认这个判断是错误的——这是纯粹的工程思维。从科研角度看，正是因为四旋翼在工程上的很多问题已经解决，科研工作者不需要从头搭建硬件和底层控制，可以直接在成熟平台上验证高层算法，反而使其成为极好的科研平台。这一反思揭示了工程思维与科学思维的根本差异。
\end{knowledgebox}

\subsection{传统控制理论在RL时代的价值}

杨硕讲述了与南方科技大学张威教授的一次重要对话。张威教授是传统控制理论背景出身，教授 Lyapunov 稳定性、状态空间模型等硬核数学课程，但后来迅速转向强化学习。

杨硕问他："你会不会觉得自己过去这么多年的积累都给浪费掉了？"张威的回答是否定的。他认为\textbf{强化学习本质上是一个非常 model-based 的东西}，理由如下：

\begin{itemize}
    \item 当前强化学习的许多进展来自 NVIDIA 将物理仿真搬到GPU上
    \item 物理仿真本质上是极端 model-based 的方法，基于系统运动学建模和优化求解
    \item 仿真中涉及的 Euler method、estimation、contact constraint 处理等，与传统MPC和机器人学高度一致
    \item 不管强化学习算法如何演进，仿真接触力时仍然需要将复杂约束转换为 linear complementarity constraint 然后求解——这些基本物理规律不会改变
    \item 人们对模型的理解被 encode 在仿真器中，而神经网络只是在 implicitly 学习仿真器中的物理知识
\end{itemize}

\begin{importantbox}{传统控制知识在强化学习中的延续}
张威教授的观点给出了一个深刻洞察：强化学习并非对传统控制的否定，而是在更高层面的融合。控制理论、系统建模的知识被编码进物理仿真器，神经网络通过强化学习隐式地学习了这些知识。因此，下一代机器人研究者仍然需要学习信号系统等基础知识，但不必像前辈那样穷尽 Kalman Filter 的四种推导方式，而应更多关注这些知识与 fundamental 数学的连接，以及如何在神经网络体系中将传感器数据和图像数据抽象为低维结构以辅助机器人执行任务。
\end{importantbox}

\subsection{科研到产品的距离}

杨硕认为科研工作与实际产品之间的距离"毫无疑问非常大"，但大的部分不在技术本身，而在于技术进入产品后如何 fit in 完整的产品生命周期——从技术实现到量产、售后、返修。每一个环节都至关重要，都有大量细碎的问题需要解决。

他特别指出，对于硬件产品来说，每一个环节都可能涉及技术，甚至可能蕴含科学问题，但解决起来更多是一个系统化的过程。你必须有返修售后相关的人，收集数据做分析——说起来所有公司都要做这些事情，但实际去做的时候"会有很多很细碎的问题需要去关注"。这种系统化的执行能力，往往是学术界训练中缺失的环节。

\begin{warningbox}{科研到产品的距离不在技术本身}
许多PhD学生倾向于认为"技术做出来就差不多了"，但杨硕的大疆和Tesla经验揭示了一个反直觉的事实：技术实现可能只是产品化的第一步。量产的良率控制、供应链管理、售后服务体系、返修流程设计——每一环都同样重要且充满挑战。一个技术demo到一个可以大规模交付并长期使用的产品之间，存在着从"说起来都知道"到"做起来全是坑"的巨大鸿沟。
\end{warningbox}

\subsection{本章小结}

CMU的博士经历让杨硕完成了从工程思维到科学思维的转变。冯·卡门关于科学家与工程师的区分、张威教授关于传统控制在RL时代延续价值的洞察，以及对科研与产品距离的清醒认识，共同塑造了他兼具学术深度与工程务实的独特视角。


